\section{Experimental Setup}
\label{sec:setup}

\paragraph{Datasets.}
We use \V{nDatasets} classification datasets (\cref{tab:datasets} in the appendix).
The \V{nOther} FG datasets cover birds, aircraft, cars, dogs, flowers, food, pets, vegetables and fruit, anime characters, and natural species~\cite{wah2011cub,vanhorn2015nabirds,maji2013fine,krause2013cars,khosla2011dogs,nilsback2008flowers,bossard2014food,parkhi2012pets,hou2017vegfru,vanhorn2018inat} \TBD{citation for Moe}.
The \V{nSc} UFG datasets are the four soybean cultivar datasets and the cotton cultivar dataset of the ultra-fine-grained benchmark~\cite{yu2021benchmark}.
The grouping follows the benchmark the data come from, not a measured property.
Averaged over datasets, baseline top-1 is \V{fpAccSc}\% for the UFG group and \V{fpAccOther}\% for the FG group.

\paragraph{Backbones, fine-tuning strategies and quantization.}
We use ViT-B/16~\cite{dosovitskiy2021vit}, Swin-B~\cite{liu2021swin} and BEiTv2-B~\cite{peng2022beitv2}, each trained on every dataset with three fine-tuning strategies.
Full fine-tuning (FT) updates the whole network, linear probing (LP) trains only the classifier on a frozen backbone, and CAL fine-tunes with counterfactual attention learning~\cite{rao2021cal} \TBD{confirm LP and the pretrained weights}.
A backbone and a strategy form one setting, and each dataset in one setting defines one context, giving 144 contexts.
HQQ~\cite{badri2023hqq} quantizes the weights of every linear layer of the backbone and keeps the classifier at full precision; activations are not quantized \TBD{HQQ version}.
The grid covers bit widths $b \in \{1, 1.58, 2, 3, 4, 5, 6, 8\}$ and group sizes $g \in \{8, 32, 128, 512\}$, giving \V{nConditions} conditions, written W3/G128 for $b=3$ and $g=128$.
The baseline is the same checkpoint evaluated without HQQ under FP16 mixed precision on the same test images.
The grid holds \V{cells} results, one for every condition of the \V{contexts} contexts.

\begin{figure*}[t]
  \centering
  \includegraphics[width=0.92\linewidth]{figures/gap_heatmaps.pdf}
  \caption{Retention gap (FG $-$ UFG) for each backbone and fine-tuning strategy, by bit width and group size. Positive values mean the UFG datasets retain less. Group means weight the datasets equally.}
  \label{fig:gap}
\end{figure*}

\paragraph{Retention and difficulty.}
For baseline accuracy $a_{\mathrm{FP}}$ and quantized accuracy $a_{\mathrm{Q}}$, \term{accuracy retention} is $r = a_{\mathrm{Q}} / a_{\mathrm{FP}}$, reported as accuracy recovery in LLM quantization~\cite{kurtic2025bf16}, and relative loss is $1-r$.
Group retention weights the datasets of a group equally, and the \term{retention gap} is the FG group's retention minus the UFG group's.
The primary conditions are W3 and W4 at group sizes 8 and 128, where retention is neither saturated near zero nor near one.
Within each setting, the \term{difficulty score} of a dataset is its percentile rank of relative loss averaged over the four primary conditions.
Kendall's $W$ measures how well the four conditions agree; we treat the score as a common ranking only when $W \ge 0.6$.

\paragraph{Decision margin.}
For baseline logits $\bm{z}$ of an image $x$, the \term{normalized decision margin} is
\begin{equation}
  m(x) = \bigl(z_{(1)} - z_{(2)}\bigr) / \sigma(\bm{z}),
  \label{eq:margin}
\end{equation}
where $z_{(1)}$ and $z_{(2)}$ are the two largest logits and $\sigma(\bm{z})$ is their standard deviation over classes.
The margin measures the distance to the second-ranked class only, while quantization can also make a lower-ranked class the winner.
The \term{vulnerable-tail margin} $q_{10}$ is the 10th percentile of $m(x)$ over the test images of a context.
We prespecify $q_{10}$ as the primary margin predictor.
It gives a difficulty signal when its Spearman correlation with the difficulty score is $\rho \le -0.5$ with a one-sided dataset permutation $p < 0.05$.
The correlation must also be negative in at least three of the four primary conditions and after leaving out any one dataset.
We further test whether adding $q_{10}$ to the prespecified controls, baseline accuracy and the log number of classes, lowers the leave-one-dataset-out error of a linear model of the score.

\paragraph{Margin matching at W3/G128.}
To ask how much of the gap follows from the different margin distributions, we reweight the test images of both groups to their common margin distribution.
With bins of width 0.05 and datasets weighted equally within each group, the target is the bin-wise minimum of the two group histograms.
Group retention becomes weighted quantized accuracy over weighted baseline accuracy.
We report the overlap of the histograms and the effective sample size (ESS) beside the matched gap.
The feature and decision diagnostics at W3/G128 are defined in \cref{sec:diagnostics}.
